% ============================================================
%  Paper A — Extended Data Figures
%  Figures moved from main text to fit Nature Plants' 6-display-item
%  limit on Article submissions. Each figure keeps its original
%  label (fig:FA2, fig:M3, etc.) so existing cross-references
%  resolve unchanged. Numbering is ED1, ED2, ... via \renewcommand.
% ============================================================

\setcounter{figure}{0}
\renewcommand{\figurename}{Extended Data Figure}
\renewcommand{\thefigure}{\arabic{figure}}

% --- Moved figures (in recommended reading order) ---
% Captions will auto-prefix "Extended Data Figure N:"; cross-references
% in the main text use the form "Extended Data Fig.~\ref{fig:XXX}"
% which resolves to "Extended Data Fig. N" (no double-"ED").

% FA2: Model organism decline and crop species ascent
\begin{figure}[H]
  \centering
  \includegraphics[width=0.72\textwidth]{FA2_model_organisms}
  \caption{\textbf{Model organism decline and crop species ascent.}
    Relative publication share (as percentage of annual plant science output)
    for \textit{Arabidopsis thaliana} (red), rice (blue), wheat (green), and
    maize (orange), 1990--2024.
    Vertical dashed line marks 2012, the year of peak
    \textit{Arabidopsis} relative share (see Discussion for the
    qualitative-nexus interpretation; the logistic growth inflection
    at 2006 is shown in Fig.~\ref{fig:FA1}).
    Secondary $y$-axis (right) shows absolute annual publication counts for
    \textit{Arabidopsis} to illustrate that relative decline precedes
    absolute decline.}
  \label{fig:FA2}
\end{figure}

% M2: Arabidopsis versus crops head-to-head share comparison
\begin{figure}[H]
  \centering
  \includegraphics[width=0.72\textwidth]{M2_arabidopsis_vs_crops}
  \caption{\textbf{Arabidopsis versus crop species: head-to-head share
    comparison.}
    Direct comparison of \textit{Arabidopsis} (red) against the combined
    rice + wheat + maize share (blue), normalised to organism-specific papers
    only.
    The crossover point marks the transition from model-organism to crop-centred
    output majority.}
  \label{fig:M2}
\end{figure}

% M3: Applied versus fundamental paradigm by organism
\begin{figure}[H]
  \centering
  \includegraphics[width=0.72\textwidth]{M3_paradigm_shift}
  \caption{\textbf{Applied versus fundamental paradigm by organism.}
    Stacked bar chart of paradigm classification for each organism category.
    The field-wide applied share is 80\%; \textit{Arabidopsis} (95.6\%) is the
    most applied organism.
    Chi-square test: $\chi^2 = 12{,}361$, $df = 9$, $p \approx 0$.}
  \label{fig:M3}
\end{figure}

% M8: CRISPR-Cas9 adoption curves by organism
\begin{figure}[H]
  \centering
  \includegraphics[width=0.72\textwidth]{M8_crispr_adoption}
  \caption{\textbf{CRISPR-Cas9 adoption curves by organism.}
    Annual CRISPR publication share for each organism category, 2012--2024.
    \textit{Arabidopsis} (red), rice (blue), and wheat (green) lead adoption;
    maize (orange) lags by approximately one year.
    Fitted logistic S-curves are shown as solid lines.}
  \label{fig:M8}
\end{figure}

% M4: Topic landscape — 30 macro-themes from BERTopic
\begin{figure}[H]
  \centering
  \includegraphics[width=0.52\textwidth]{M4_topic_landscape}
  \caption{\textbf{Topic landscape of plant science (SPECTER2 + UMAP).}
    UMAP projection of SPECTER2 embeddings for a sample of $n = 49{,}446$
    classified papers, coloured by organism category
    (\textit{Arabidopsis}, barley, cotton, maize, rice, soybean, tomato,
    tobacco, potato, wheat, other crop, other model organism; see
    legend). The 30 macro-themes obtained from BERTopic modelling of
    the full 1{,}632 fine-grained topics are described in the main
    text and visualised by their temporal trajectories in
    Fig.~\ref{fig:M5}.}
  \label{fig:M4}
\end{figure}

% M5: Topic evolution heatmap (1990–2024)
\begin{figure}[H]
  \centering
  \includegraphics[width=0.60\textwidth]{M5_topic_evolution_heatmap}
  \caption{\textbf{Macro-theme growth/decline by decade (top 20 most
    variable macro-themes).}
    Heatmap of decadal growth rate for the 20 macro-themes with
    largest share change across 1990s--2020s.
    Cell colour gives the percentage-point change in publication
    share per decade (red = gaining share, blue = losing share).
    Row labels are the top BERTopic terms for each macro-theme
    (space-limited raw tokens; human-interpretable theme names given
    in main text). Across the 30 macro-themes (not all shown),
    Shannon entropy of the distribution has declined from
    1.82 bits (1990s mean) to 1.14 bits (2020s mean), indicating
    concentration of the research agenda.}
  \label{fig:M5}
\end{figure}

% M6: Organism-topic enrichment matrix
\begin{figure}[H]
  \centering
  \includegraphics[width=0.72\textwidth]{M6_organism_topic_enrichment}
  \caption{\textbf{Organism-topic enrichment matrix.}
    Log-odds enrichment of each organism category for each macro-theme,
    relative to the field baseline.
    Significant enrichments ($p < 0.01$ after Bonferroni correction) are
    indicated by filled cells.}
  \label{fig:M6}
\end{figure}

% M7: Semantic drift — cosine distance of organism centroids from Arabidopsis
\begin{figure}[H]
  \centering
  \includegraphics[width=0.72\textwidth]{M7_semantic_drift}
  \caption{\textbf{Semantic drift: cosine distance of organism centroids
    from \textit{Arabidopsis} (SPECTER2 embeddings).}
    Lower values indicate greater conceptual similarity to
    \textit{Arabidopsis}.
    Cotton (olive) and rice (blue) show the strongest convergence over
    2010--2024.
    Shaded band: 95\% bootstrap confidence interval.}
  \label{fig:M7}
\end{figure}

% FA5: Six intellectual epochs defined by concept marriages
\begin{figure}[H]
  \centering
  \includegraphics[width=0.60\textwidth]{FA5_concept_marriages}
  \caption{\textbf{Six intellectual epochs defined by concept marriages.}
    (\textbf{A}) Annual count of newly emerging concept marriages
    (new keyword pairs whose co-occurrence rose above the
    significance threshold) in the plant-science knowledge space,
    1960--2024. Background shading distinguishes the six epochs
    identified by hierarchical clustering of the concept
    co-occurrence matrix: 1960--1974, 1975--1984, 1985--1994,
    1995--2004, 2005--2014, 2015--2024.
    (\textbf{B}) Total concept marriages per epoch, showing a monotonic
    decline from 1{,}123 (1960--1974 cohort) to 163 (2015--2024
    cohort), consistent with a maturing knowledge space in which
    fewer genuinely novel keyword combinations emerge each year.
    The 2015--2024 epoch is dominated by crop-genomics and
    precision-breeding marriages (main text).}
  \label{fig:FA5}
\end{figure}

% End of Extended Data; subsequent SI figures will continue the ED counter,
% which is fine because SI figures use manually-labeled captions (see
% supplementary.tex "Supplementary Figure S1..." headings that don't rely
% on auto-numbering).
